\section{How to get involved in the development of learning mechanics}
\label{sec:how_to_get_involved}

It is always difficult to start doing research in a new field.
Consequently, we would like to make it as easy as possible for newcomers to get started.
In this section, we extend a hand with some encouragement and advice.

There is no specific academic background required to do useful work in this field.
Well-regarded researchers in deep learning theory come from backgrounds in physics, mathematics, computer science, neuroscience, statistics, and more.
Moreover, knowing another field well is useful, and established ideas from other fields can be applied to deep learning in some form or another, as the diversity of perspectives on deep learning theory attests (\cref{sec:learning_mech_and_mechinterp}).
Good things grow from cross-pollination.
A firm grasp of undergraduate mathematics, a familiarity with deep learning, and a desire to learn are the only definite prerequisites.

\subsection{Tenets for getting started}
If you want to join this field, you are more than welcome.
While there is no single correct way to craft theory, there are plenty of pitfalls that many of us encountered when starting out.
To help avoid some of them, we have compiled a shortlist of guiding principles for doing research in this field.
These tenets are not intended to maximize your number of citations in the short term, and following them may involve some swimming against the current of academia.
Instead, they are intended to maximize your impact in the long term and your ability to integrate and contribute to the community.

\begin{enumerate}

    \item
    \textbf{Do experiments frequently.}
    As discussed in \cref{sec:reason-measurable}, deep learning is a field where the cost of doing experiments is relatively low, with a fast turnaround time.
    Use that to your advantage!
    Experiments serve to check assumptions, inform theoretical models, reveal the limitations of a theory, peer beyond the cases that a theory covers, and surface interesting phenomena to study in the first place.
    Try to include experiments in every paper, and make them as simple and revealing as possible.

    \item
    \textbf{Simplicity and insight matter more than technical complexity.}
    If you want to do work that is useful for others, they need to understand it, not merely be impressed by it.
    Take the time to simplify your findings, identify the underlying intuition, and check with simple experiments.
    A useful idea for thinking about a \textit{type} of problem is generally more valuable than a difficult theorem or a solution to any \textit{particular} problem, so emphasize these useful ideas when you present your work.
    This will make your results more accessible and easier for others to extend and apply.
    Your conference reviewers may disagree with this philosophy --- the conference system tends to reward technical complexity and undervalue simplicity --- but it will make your work more impactful in the long run.


    \item
    \textbf{Value scientific understanding over state-of-the-art performance.}
    Applied deep learning is a field of engineering whose progress is measured by benchmarks.
    For scientists of deep learning, the game is different: your contribution is gauged by your contribution to collective understanding.
    It is easy to feel pressure to tack on some engineering benefit in a scientific paper to make the paper seem more relevant and timely, but doing so usually dilutes the paper's scientific contribution without really affecting practice.
    Of course, fundamental science should eventually improve state-of-the-art performance, and when this naturally falls out of the science, it is a powerful way to demonstrate what has been understood.%
    \footnote{For a good example of this, see \citet{yang2022tensor}.}
    When it does not, though, there is no need to force it: set benchmarks aside for the time being and seek understanding on its own terms.

    \item
    \textbf{Don't try to do it alone.}
    Deep learning is a field with a lot of history and many known results, and guidance from a live human will help you.
    You can get pretty far from reading material online and talking to AI, but it doesn't replace human collaboration and mentorship.
    You should seek out other people interested in this area, ask for their feedback, and ask to work with them.
    A weak corollary to this is: when one exists, watch the talk version of a paper in addition to reading it.
    A great deal more of the nuance of a project is conveyed through in a live presentation.
    The \href{http://physicsmeetsml.org/}{Physics Meets ML} and \href{https://www.physicsoflearning.org/webinar-series}{Physics of Learning} series are good places to look for talks on what we here call learning mechanics.

    \item
    \textbf{Try a few different problems before going deep into one.}
    Deep learning theory has so many open questions that you probably won't know where to start.
    That's okay --- just jump in, and feel free to change problems a few times early on.%
    \footnote{Most authors of this paper did this, as you can see from our respective research records!}
    Knowing multiple areas is essential to having high-level ideas anyways, and working in an area is the best way to learn it.

    \item
    \textbf{Invest in fundamental tools and techniques.}
    Compared to more established fields of science, mathematics, and engineering, deep learning is very young.
    These other fields have identified deep ideas and developed powerful tools applicable to broad classes of problem.
    Learning these fundamental tools comes in handy when similar problems arise in deep learning.
    For example, statistical physics and random matrix theory have powerful tools for thinking about high-dimensional interacting systems, and these tools are of great use in learning mechanics when one takes infinite limits.
    Many of the central ideas of classical optimization theory make regular appearances in the study of neural network optimization.
    Other concepts from statistical signal processing, such as wavelet decompositions, graphical models, and information theory, have been comparatively less leveraged so far, but may prove to be useful and complementary tools of learning mechanics.


\end{enumerate}


We will put as much useful introductory material as we can on \websiteurl, and we encourage discussion in the comments there.
We also encourage taking a crack at the open directions in \cref{sec:open_dirs}.
Work hard, have fun, and best of luck --- we hope to see a great deal more fundamental science of deep learning in the next few years!
